\documentclass[11pt,a4paper]{article}
\usepackage[T1]{fontenc}
\usepackage{lmodern,amsmath,amssymb,amsthm,mathtools}
\usepackage[margin=28mm]{geometry}
\usepackage{microtype,booktabs,array}
\usepackage[colorlinks=true,linkcolor=blue,citecolor=blue,urlcolor=blue]{hyperref}
\usepackage{xurl}
\hypersetup{pdftitle={The exact maximum in Erdos problem 272 for sufficiently large ground sets},
 pdfauthor={Li Xiang},pdfsubject={Computer-assisted partial solution of Erdos problem 272}}
\newtheorem{theorem}{Theorem}
\newtheorem{proposition}[theorem]{Proposition}
\newtheorem{lemma}[theorem]{Lemma}
\newtheorem{corollary}[theorem]{Corollary}
\theoremstyle{definition}\newtheorem{definition}[theorem]{Definition}
\newcommand{\F}{\mathcal F}
\newcommand{\G}{\mathcal G}
\newcommand{\D}{\mathcal D}
\newcommand{\A}{\mathcal A}
\newcommand{\Hh}{\mathcal H}
\newcommand{\repo}{\url{https://github.com/lixiang90/Erdos272}}
\newcommand{\code}[1]{{\small\nolinkurl{#1}}}
\setlength{\emergencystretch}{2em}
\title{The exact maximum in Erd\H{o}s problem 272\protect\\for sufficiently large ground sets}
\author{Li Xiang\thanks{Project and complete proof sources: \repo.}}
\date{30 October 2026}
\begin{document}
\maketitle
\begin{abstract}
Let $v(N)$ be the maximum cardinality of a family of distinct subsets of
$\{1,\ldots,N\}$ whose intersections of distinct members are nonempty
arithmetic progressions. We give a computer-assisted proof that
\[
 v(N)=\binom N2+1+\left\lfloor\frac{N-1}{4}\right\rfloor
 \qquad(N\ge 2^{49}).
\]
The lower bound is Szab\'o's construction. For the upper bound we delete
large nonprogressions, obtain a common point for the remaining
nonprogressions, and force all members avoiding that point to be long
intervals. Their common interval then pays for the deleted members.
We give the explicit numerical estimates and the precise interfaces of
the formally proved counting lemmas. The accompanying Lean development is
part of the proof. More refined computations in that development establish
the same formula for every $N\ge642498$; those refinements are not developed
in this paper. The remaining finite range is not settled here.
\end{abstract}

\section{Statement and scope}
Write $[N]=\{1,\ldots,N\}$. An arithmetic progression (AP) means a
nonempty finite set $\{a,a+d,\ldots,a+kd\}$ with $d\ge1$ and $k\ge0$.
Thus singletons and pairs are APs. A family $\F\subseteq2^{[N]}$ is
\emph{legal} if $S\cap T$ is an AP whenever $S,T\in\F$ and $S\ne T$.
The members themselves need not be APs. Set
\[
 B(N)=\binom N2+1+\left\lfloor\frac{N-1}{4}\right\rfloor,
 \qquad v(N)=\max_{\F\text{ legal}}|\F|.
\]
The maximum exists because its domain is finite.

\begin{theorem}\label{thm:main}
For every integer $N\ge2^{49}$, $v(N)=B(N)$. For every $N\ge1$,
there is a legal family of size $B(N)$.
\end{theorem}

This is a partial answer to Erd\H{o}s problem 272 and the corresponding
Justin Sun Prize problem JSP-000233. Szab\'o established the leading term
and the construction giving $B(N)$~\cite{Szabo}. Yang formulated the exact
formula and proved the common-point bound using private pairs and a
staircase inequality~\cite{Yang}. Our formal development also follows the
non-common-point structural route of Kohlmeyer and Kruer~\cite{KK}.
We do not claim priority for those constructions or methods. We prove
existence of an optimal common-point family, not that every optimal family
has a common point.

The theorem is \emph{computer-assisted}: the finite-set counting lemmas
specified below have complete proofs in the companion Lean source, which
forms an electronic supplement to this paper. The paper supplies their
mathematical interfaces and the full numerical and logical assembly.
It does not replace those source proofs by an assertion that a numerical
search succeeded. Section~\ref{sec:verification} gives direct paths and
reproduction commands.

\section{The matching construction}
Let $c=\lceil N/2\rceil$ and $k=\lfloor(N-1)/4\rfloor$. Start with every
subset of $[N]$ of size at most three containing $c$. For each $1\le d\le k$,
remove the two triples
\[
 \{c-2d,c,c+d\},\qquad \{c-d,c,c+2d\},
\]
and add the five-term progression
$P_d=\{c-2d,c-d,c,c+d,c+2d\}$ and its two four-term end deletions.

\begin{proposition}[Szab\'o's construction]\label{prop:lower}
For every integer $N\ge1$, the resulting family is legal and has cardinality $B(N)$.
\end{proposition}
\begin{proof}
Since $2k\le\min(c-1,N-c)$, all its points lie in $[N]$.
Every member contains $c$. Intersections involving a member of size at
most two, or two distinct triples, therefore have size one or two.
Two APs have an AP as their nonempty intersection: their common points
form a residue class modulo the least common multiple of their steps,
restricted to their common span.

It remains to intersect a retained triple with an added AP. Only a
three-point intersection needs checking. Such a triple lies in some
$P_d$ and contains $c$. After translating and scaling, the six possible
triples through zero in $\{-2,-1,0,1,2\}$ consist of four APs and the
two deleted triples $\{-2,0,1\}$ and $\{-1,0,2\}$.
Thus every required intersection is an AP.

The ordered gaps $(2d,d)$ and $(d,2d)$ distinguish the deleted triples.
The step and endpoints distinguish the added APs. Consequently exactly
$2k$ triples are removed and $3k$ larger members are added, giving
\[
 N+\binom{N-1}{2}-2k+3k=\binom N2+1+k.
\]
This also covers $k=0$.
\end{proof}

\section{The structural counting certificates}\label{sec:structural}
A member is \emph{crooked} if it is not an AP. For a legal family $\F$,
define
\[
 \D(\F)=\{S\in\F:S\text{ is crooked and }28|S|>N\},
 \qquad \G(\F)=\F\setminus\D(\F).
\]
Every AP member is retained. Write
$H_N=\sum_{i=1}^N1/i$ and $J=\lfloor\log_2N\rfloor+1$.
All intervals below are integer intervals, with both endpoints included.

We use the following formally proved structural facts. Their hypotheses
refer to the actual family and its subsets. In particular they do not
assume an upper bound on an abstract collection of witnesses.

\begin{lemma}[Preparation and endpoint deficit]\label{lem:prepare}
Suppose $N\ge2^{49}$, $\F$ is legal, and $|\F|>B(N)$.
Either leave $\F$ unchanged or reflect every member by $x\mapsto N+1-x$,
obtaining a legal family $\Hh$ with $|\Hh|=|\F|$.
There are integers $1\le c<L\le R\le N$ with the following properties.
Put $\D=\D(\Hh)$, $\G=\G(\Hh)$, and
$\A=\{S\in\G:c\notin S\}$.
\begin{enumerate}
\item Every crooked member of $\G$ contains $c$, and $\A\ne\varnothing$.
\item Every $S\in\A$ has $\min S\le L$ and $\max S\ge R$.
\item With $t=L+N-R$, there is a nonnegative witness count $E$ such that
\begin{align}
 |\D|&\le3137N,\qquad E\le6NH_N^3,\label{eq:deletion}\\
 |\Hh|&\le\frac{N^2}{2}-\frac{t^2}{102}
  +|\D|+E+40NH_N^2+6NH_N+24N+(L-c)+3.
  \label{eq:deficit}
\end{align}
\end{enumerate}
\end{lemma}

Here and below the statements of the source lemmas are stronger than the
specializations displayed. For Lemma~\ref{lem:prepare}, the exact source
is \code{Counting/MiddleDeficit52.lean}, declaration
\code{actual_original_prepared_endpoint102}; deletion and the exceptional
count are supplied by \code{Core/DenseSharper41.lean} and
\code{Arithmetic/DivisorMoment45.lean}.
All source paths in the paper are relative to
\code{formal/Erdos272/}.

We describe the counting mechanism to clarify what is certified. A crooked
set containing a root $c$ has a pair $\{x,y\}$ for which the smallest AP
containing $c,x,y$ is not contained in the set. Such a pair cannot belong
to another member through $c$, since their legal intersection would
contain that AP. These private pairs, and the pairs covered by AP
members, are counted in the same pair universe. Deleting large crooked
members first bounds their cost by $(4\cdot28^2+1)N=3137N$.
The remaining crooked members have a common point by the empty-kernel
counting bound. If the retained family itself has a common point, the
common-point bound and the omitter estimate already exclude a
counterexample. Otherwise an avoiding AP supplies a one-sided
orientation, and endpoint extrema supply $L,R$.

The endpoint estimate counts primitive pairs in the middle region;
coprimality, floor errors, and exceptional witnesses are retained in
\eqref{eq:deficit}. The coefficient $1/102$ is deliberately conservative.
These combinatorial arguments, including the common-point bound and its
finite staircase certificate and unbounded tail, are proved in the import
closure of the cited module. No kernel conjecture is assumed.

\begin{lemma}[Excluding noninterval avoiders]\label{lem:noninterval}
Let $\G$ be legal, let $c\in[N]$, and suppose every crooked member of
$\G$ contains $c$. Suppose $c\le L\le N$, and each member avoiding $c$
contains a point at most $L$. If an avoiding member is not an integer
interval, then
\begin{equation}\label{eq:noninterval}
 |\G|\le\frac{3N^2}{8}+\frac53LN+2NH_N+20N+3+6NH_N^3.
\end{equation}
In particular, if $N\ge6$, $N<2^J$, $J\ge2$, $32L\le N$, and
$16384J^3\le N$, then $|\G|\le11N^2/24$.
\end{lemma}

The exact finite-set counting proof is
\code{Global/NonInterval53.lean}, declarations
\code{actual_noninterval_family_count} and
\code{actual_noninterval_family_gap}.
An avoiding member is an AP. If it is not an interval, its step is at
least two and its size is at most $N/2+1$. The unordered pairs meeting
it number at most $3N^2/8+N$. Private witnesses for crooked members
inject into these pairs, except for the separately bounded exceptional
members. Counting the AP members with a low endpoint gives
\eqref{eq:noninterval}. For the stated consequence, use $H_N\le J$
and $J\ge2$. The last four terms in \eqref{eq:noninterval} are at most
$31NJ^3\le31N^2/16384$. Hence
\[
 \frac38+\frac5{96}+\frac{31}{16384}<\frac{11}{24}.
\]
The elementary inequality $H_N\le J$ follows by partitioning $[N]$
into the dyadic blocks $\{1\},[2,3],[4,7],\ldots$; each contributes
at most one.

\begin{lemma}[Interval count and restoration]\label{lem:restore}
Let $N\ge29^3$ and $1\le c<L\le R\le N$, and put $t=L+N-R$.
Let $\Hh$ be legal, $\D=\D(\Hh)$, and $\G=\G(\Hh)$.
Suppose the members of $\G$ avoiding $c$ form a nonempty family of
intervals $[a,b]$ satisfying $c<a\le L\le R\le b\le N$.
If
\begin{equation}\label{eq:payments}
 6tJ\le N,\qquad812t\le N,
\end{equation}
then
\begin{equation}\label{eq:restore-counts}
 |\G|\le\binom N2+1,
 \qquad |\D|\le131(t-1)\le\left\lfloor\frac{N-1}{4}\right\rfloor.
\end{equation}
Consequently $|\Hh|\le B(N)$.
\end{lemma}

The first bound is
\code{Retention/Terminal36.lean}, declaration
\code{actual_native_interval_terminal}.
The second is \code{Retention/JointFringe44.lean}, declaration
\code{actual_dense_fringe_bound}, with parameters $r=28$, $p=29$.
Their assembly is \code{Global/GlobalUpper53.lean}, declaration
\code{actual_deleted_long_core_restore28}.

Here is why the two counts fit together. An avoiding interval contains
$[L,R]$, so this core lies in an actual AP member of the original family.
Its exterior in $[N]$ has cardinality $t-1$. A deleted member has more
than $N/28$ points. The inequality $28\cdot29t\le N$ forces a long
trace in the core. Bounded-step trace anchors and points outside the core
give the joint bound
\[
 \operatorname{JointCount}(29)(t-1)
 =\left(29+\sum_{d=1}^{28}\left\lfloor\frac{29}{d}\right\rfloor\right)(t-1)
 =131(t-1).
\]
The aligned and off-residue cases are both included in this count.
For the final payment, $4\cdot131t=524t\le812t\le N$ implies
$4\cdot131(t-1)\le N-1$. Since $\D$ and $\G$ partition the original
family, their bounds add to exactly $B(N)$.

\section{The explicit large-order argument}
We now prove Theorem~\ref{thm:main} from the certified counting lemmas.
Every numerical estimate in this section is elementary.

\begin{lemma}[Uniform logarithmic payments]\label{lem:gates}
If $N\ge2^{49}$ and $J=\lfloor\log_2N\rfloor+1$, then
\begin{equation}\label{eq:gates}
 J\ge50,\qquad470771616J^3\le N,\qquad25704J^5\le N.
\end{equation}
\end{lemma}
\begin{proof}
Put $k=J-1\ge49$. At $k=49$, direct integer arithmetic gives
\[
 470771616\cdot50^3=58846452000000<2^{49},\qquad
 25704\cdot50^5=8032500000000<2^{49}.
\]
For $k\ge49$, $(k+2)/(k+1)\le51/50$, and $(51/50)^5<2$.
Thus both $(k+1)^3/2^k$ and $(k+1)^5/2^k$ decrease by induction.
The desired bounds follow from $2^k\le N$.
\end{proof}

\begin{proof}[Proof of Theorem~\ref{thm:main}]
Proposition~\ref{prop:lower} proves the lower bound. Suppose for a
contradiction that $N\ge2^{49}$ and $|\F|>B(N)$ for a legal family.
Apply Lemma~\ref{lem:prepare} and use its notation. Since $H_N\le J$,
$L-c\le N$, and $N\ge1$, \eqref{eq:deletion}--\eqref{eq:deficit}, together
with $|\Hh|\ge\binom N2$, imply
\begin{equation}\label{eq:t-square}
 t^2\le102N(6J^3+40J^2+6J+3166)\le714NJ^3.
\end{equation}
For the last inequality, $J\ge50$ and the identity
\[
 J^3-40J^2-6J-3166
 =(J-50)(J^2+10J+494)+21534\ge0
\]
suffice. The slack in $3166$ also covers the $N/2$ difference between
$N^2/2$ and $\binom N2$, and the constant $3$ in \eqref{eq:deficit}.

As $812^2\cdot714=470771616$ and $36\cdot714=25704$,
Lemma~\ref{lem:gates} and \eqref{eq:t-square} give
\[
 (812t)^2\le N^2,\qquad(6tJ)^2\le N^2.
\]
All quantities are nonnegative, so $812t\le N$ and $6tJ\le N$.
Also $L\le t$, whence $32L\le N$, and
$16384J^3\le470771616J^3\le N$.

If an avoiding member of $\G$ were not an interval,
Lemma~\ref{lem:noninterval} would give $|\G|\le11N^2/24$.
On the other hand, $|\Hh|>B(N)$ and $|\D|\le3137N$ imply
\[
 \frac{N^2}{2}\le |\G|+3138N.
\]
These inequalities would force $N\le24\cdot3138=75312$, contrary to
$N\ge2^{49}$. Therefore every avoiding member is an integer interval
$[a,b]$. The endpoint conditions in Lemma~\ref{lem:prepare} give
$a\le L\le R\le b$. Since $c<L$ and $c\notin[a,b]$, necessarily
$c<a$.

All hypotheses of Lemma~\ref{lem:restore} now hold, for the same
family, root, and endpoints. It follows that
$|\F|=|\Hh|\le B(N)$, a contradiction.
\end{proof}

\section{Formal verification and the stronger range}\label{sec:verification}
The repository is \repo. The public entry is
\code{formal/PublicResults.lean}. It defines the ground set, AP predicate,
legality, $B(N)$, the finite set of legal families, and its extremal value
directly in that file using Lean and mathlib types. The proofs apply to
these definitions by definitional equality with the internal finite-set
model. In particular there is no additional common-point hypothesis in
the public upper bound.

\begin{corollary}[Refined formal result]\label{cor:refined}
For every integer $N\ge642498$, $v(N)=B(N)$, and a legal family attaining
this value exists.
\end{corollary}

This stronger result uses more refined arithmetic and finite certificates
than the argument developed above. It is not proved in detail in this
paper; readers can inspect and check the complete source through
\code{PublicResults.exact_value} and \code{PublicResults.optimal_family}
in namespace \code{JSP233}. The paper's own theorem has the separate
entry \code{PublicResults.paper_exact_value}, which calls the
$2^{49}$ argument directly. Thus the simpler theorem can be traced
without relying on the refined threshold.

The development is pinned to Lean 4.35.0 and mathlib's matching stable
release. The lock file records the exact dependency revisions. From a
checkout with the specified Lean toolchain, run
\begin{verbatim}
cd formal
lake exe cache get
lake build
\end{verbatim}
The public entry prints the theorem types and their axiom dependencies.
The proofs use only \code{propext}, \code{Classical.choice}, and
\code{Quot.sound}; finite certificates are checked by Lean's kernel.
A full build can require substantial time and memory. With limited
resources, compile the project modules sequentially in dependency order,
reuse compiled dependencies, and choose process, thread, memory, and CPU
limits suited to the available hardware.

The directory groups separate foundational set geometry, private-pair
counting, retention and restoration, arithmetic certificates, and global
assembly. A proof covering the entire remaining finite range is still
missing.

\paragraph{Acknowledgements and collaboration.}
The project builds on the work cited below and was developed with
AI-assisted research, programming, and review. Formal checking and
independent review serve different purposes: the former checks the
proof terms, and the latter checks their interpretation as the stated
mathematical problem. Contributions that close the remaining finite gap
are welcome. Please credit the earlier results and all new contributions,
and coordinate a joint application for the Justin Sun Prize.

\begin{thebibliography}{9}
\bibitem{Szabo}
T.~Szab\'o, Intersection properties of subsets of integers,
\emph{European Journal of Combinatorics} \textbf{20} (1999), 429--444.
\url{https://page.mi.fu-berlin.de/szabo/PDF/aps.pdf}.

\bibitem{Yang}
Z.~Yang, Exact values and exact upper bounds for families of integers with
arithmetic progression intersections (Erd\H{o}s Problem \#272),
arXiv:2607.23004v1 (2026).
\url{https://arxiv.org/abs/2607.23004v1}.

\bibitem{KK}
J.~Kohlmeyer and L.~Kruer, \emph{A linear error bound for Erd\H{o}s problem 272},
manuscript and companion Lean development, fixed source commit
\nolinkurl{a624e9de77e0bd8868e7103f8b379f9dbd86b607} in
\href{https://github.com/JENW1N/complete-math-solutions/blob/a624e9de77e0bd8868e7103f8b379f9dbd86b607/proofs/Erdos272.lean}{JENW1N/complete-math-solutions (fixed Lean source)}.
The source informed the structural route; the present theorems are checked
in the accompanying finite-integer-set model.

\bibitem{Prize}
The Justin Sun Prize, Awards and problem list,
\url{https://github.com/TheJustinSunPrize/awards}.
\end{thebibliography}
\end{document}
